\documentclass[10pt,a4paper]{amsart}
\usepackage[margin=19mm]{geometry}
\usepackage{amsmath,amssymb,mathtools}
\usepackage[scr=rsfs,cal=euler]{mathalfa}
\usepackage{xfrac}
\usepackage[unicode,hidelinks,bookmarks=false]{hyperref}
\newcommand{\ie}{i.e.\ }
\newcommand{\cf}{cf.\ }
\newcommand{\resp}{respectively\ }
\newcommand{\Subsection}{\subsection}
\providecommand{\nobreakdash}{\nobreak\hbox{-}}
\setlength{\emergencystretch}{2em}
\multlinegap0pt
\usepackage[normalem]{ulem}
\usepackage{amssymb}
\usepackage{enumerate}
\usepackage{tikz}
\usepackage{mathtools}
\newtheorem{lemma}{Lemma}
\newtheorem{prop}{Proposition}
\newcommand{\A}{\mathcal{A}}
\newcommand{\C}{\mathcal{C}}
\renewcommand{\emptyset}{\varnothing}
\title{Iwahori Matsumoto presentation for modules of Iwahori fixed functions on symmetric spaces}
\author{Guy Shtotland}
\date{Source: arXiv:2406.16070v3 (CC BY 4.0)}
\begin{document}
\maketitle

\noindent\emph{Source and licence.} Iwahori Matsumoto presentation for modules of Iwahori fixed functions on symmetric spaces. Version arXiv:2406.16070v3 (CC BY 4.0), distributed by arXiv under the Creative Commons Attribution 4.0 International licence, http://creativecommons.org/licenses/by/4.0/, as recorded in the arXiv licence metadata for this exact version and shown on the abstract page https://arxiv.org/abs/2406.16070v3. Full licence notice: \texttt{licenses/arxiv-2406.16070-CC-BY-4.0.txt}.\par\smallskip

\noindent\textit{Editorial scope.} Counted unit: the article's prop key\_computation (source0.tex lines 934--947). The article's complete original proof of it is delivered with it (source0.tex lines 950--991, one \texttt{proof} environment of 42 lines). No mathematics is written, repaired or re-derived: the statement and the proof are the article's own lines, in the article's own notation, and no retained line is substituted or re-flowed.  Retained uncounted as the source-local context the proof needs: lines 908--912.  Nothing was omitted from any retained span, so the package carries no omission record.  No retained line contains a citation, so the wrapper carries no bibliography and no bibliography entry.  No retained line contains a figure environment, an image-inclusion command or drawing code, so no image is delivered and the package declares no asset dependency.  Editorial formatting only, in addition: an AMS \texttt{amsart} preamble that restates the article's own declarations and macros used by the retained lines, namely the environments \texttt{lemma}, \texttt{prop} on separate counters, together with the standard packages those lines need.  The retained environments print in this document as Lemma 1, Proposition 1 in order of appearance, whereas the article prints the same items with the numbering of its own full document.  The complete native source of the article as distributed in the arXiv e-print for this exact version is bundled unchanged as \texttt{sources/161229/source0.tex}.
\begin{lemma}\label{3.13}
Let $\C$ be the chamber fixed by $I$ and let $\A$ be an apartment that contains $\C$. Let $s\in \Tilde{\Delta}$ be an element that induces a reflection over a facet of $\C$. 

The set of chambers $Is\C$ is the set of all chambers that contain $\C\cap s\C$ and are different from $\C$ and its size is $q$, the size of the residue field of $F$.  
\end{lemma}

\begin{prop}\label{key_computation}

Let $s\in \Tilde{\Delta}$, $x\in H\backslash G/I$, $g$ a representative of $x$ as above. Let $f=g\C\cap gs\C$ and let $\C_1,...,\C_{q+1}$ be the chambers that contain $f$. Denote $o=\omega(g)\in\Omega$.

Let $O_{f,o}$ be the set of $H$ orbits of $(\C_1,o),...,(\C_{q+1},o)$.

Define $D_{f,o}=\sum_{x\in O_{f,o}}1_x$ and 
let $\gamma_{f,g}=\#\{1\leq i\leq q+1|H(\C_i,o)=H(g\C,o)\}$.

Then we have:
\begin{center}
    $1_x(T_s+1)=\gamma_{f,g} D_{f,o}$
\end{center}
\end{prop}

\begin{proof}

First, observe that by the definition of convolution, $1_x(T_s+1)$ is supported on $HxIsI$. And that all the $H$ orbits of chambers in $HxIs\C$ can be represented by a chamber that contains $f$.

For each $1\leq i\leq q+1$ choose $g_1,...,g_{q+1}\in G$ such that $g_i\C_0=\C_i$ and $\omega(g_i)=\omega(g)$. Denote $x_i=Hg_i$.

We now compute $1_{x}T_s(x_i)$.
Let $m$ be a Haar measure on $G$ normalized such that $m(I)=1.$
A direct calculation gives:
$$1_{x}T_s(x_i)=\int_{g'\in G} 1_{x}(x_ig'^{-1})T_s(g')dm(g')=m(\{k\in IsI|g_ik^{-1}\in HgI\}).$$ 
On the other hand, 
$$m(\{k\in IsI|g_ik^{-1}\in HgI\})=m(IsI\cap Ig^{-1}Hg_i)=m(IsIg^{-1}_i\cap Ig^{-1}H).$$
To compute the intersection, we use Lemma \ref{3.13} to write $g_iIsI=\cup_{k\neq i}g_kI$

Taking inverses, we obtain $$IsIg^{-1}_i=\cup_{k\neq i}Ig^{-1}_k$$

Passing to measures, we get

$$m(IsIg^{-1}_i\cap Ig^{-1}H)=\sum_{k\neq i}m(Ig^{-1}_k\cap Ig^{-1}H)$$

Thus we obtain
 $$1_{x}T_s(x_i)=\sum_{k\neq i}m(Ig^{-1}_k\cap Ig^{-1}H)$$


We observe that, by the our choice of $g_{1},...,g_{q+1},$ for every $k$ we either have $Ig_k^{-1}\subset Ig^{-1}H$ or $Ig_k^{-1}\cap Ig^{-1}H=\emptyset.$

So the aforementioned sum is equal to

\[
    \beta_{g\C,g_i\C}=\sum_{k\neq i}\delta_{Hg_kI,HgI} \label{eq:special} \tag{*}
\]

Here, $\delta_{Hg_kI,HgI}$ is $1$ if $Hg_kI=HgI$ and zero otherwise.

The ''missing term'' in the formula 
$$1_{x}T_s(x_i)=\sum_{k\neq i}\delta_{Hg_kI,HgI}$$ motivates us to consider the quantity 
$1_{x}T_s(x_i)+1_{x}(x_{i}).$ which is independent of $i$ and equals $\gamma_{f,g}$.

Thus we obtain $1_x(T_s+1)(x_i)=\gamma_{f,g}$.

This computation also shows that $1_x(T_s+1)$ is supported on $\{x_1,...,x_{q+1}\}$ so we obtain $1_x(T_s+1)=\gamma_{f,g} D_{f,o}$
\end{proof}

\end{document}
